% answers.tex —— 2012 年·辽宁卷·理科 · 参考答案与详细解析片段
% 由 tools/generate_projects.py 依来源总账自 2012/2012_数学真题_辽宁理_试题.md + 2012/2012_数学真题_辽宁理_解析正文.md + 2012/2012_数学真题_辽宁理_答案要点.md 生成，勿手改（改源 MD 后重跑）


\begin{answer}
\gksolhead{1}{}

\textbf{题目}：

已知全集 $U=\{0,1,2,3,4,5,6,7,8,9\}$，集合 $A=\{0,1,3,5,8\}$，集合 $B=\{2,4,5,6,8\}$，则 $(\complement_U A)\cap(\complement_U B)=$\mbox{（　　）}

（A）$\{5,8\}$；（B）$\{7,9\}$；（C）$\{0,1,3\}$；（D）$\{2,4,6\}$。

\Needspace{4\baselineskip}
\textbf{解析}

\Needspace{5\baselineskip}
先求并集：
\nopagebreak[4]
\[A\cup B=\{0,1,3,5,8\}\cup\{2,4,5,6,8\}=\{0,1,2,3,4,5,6,8\}.\]
\Needspace{5\baselineskip}
由德摩根律
\nopagebreak[4]
\[(\complement_U A)\cap(\complement_U B)=\complement_U(A\cup B).\]
\Needspace{5\baselineskip}
在 $U=\{0,1,\ldots,9\}$ 中去掉 $A\cup B$ 的元素，剩下
\nopagebreak[4]
\[\complement_U(A\cup B)=\{7,9\}.\]

选 B。

\gkanswer{答案：B。$(\complement_U A)\cap(\complement_U B)=\complement_U(A\cup B)=\{7,9\}$。}

\end{answer}
\begin{answer}
\gksolhead{2}{}

\textbf{题目}：

复数 $\dfrac{2-\mathrm i}{2+\mathrm i}=$\mbox{（　　）}

（A）$\dfrac35-\dfrac45\mathrm i$；（B）$\dfrac35+\dfrac45\mathrm i$；（C）$1-\dfrac45\mathrm i$；（D）$1+\dfrac35\mathrm i$。

\Needspace{4\baselineskip}
\textbf{解析}

\Needspace{5\baselineskip}
分子分母同乘 $2-\mathrm i$：
\nopagebreak[4]
\[\frac{2-\mathrm i}{2+\mathrm i}
=\frac{(2-\mathrm i)^2}{(2+\mathrm i)(2-\mathrm i)}
=\frac{4-4\mathrm i+\mathrm i^2}{4+1}.\]
\Needspace{5\baselineskip}
利用 $\mathrm i^2=-1$，得
\nopagebreak[4]
\[\frac{3-4\mathrm i}{5}=\frac35-\frac45\mathrm i.\]

选 A。

\gkanswer{答案：A。$\dfrac{2-\mathrm i}{2+\mathrm i}=\dfrac{(2-\mathrm i)^2}{5}=\dfrac35-\dfrac45\mathrm i$。}

\end{answer}
\begin{answer}
\gksolhead{3}{}

\textbf{题目}：

已知两个非零向量 $\boldsymbol a$，$\boldsymbol b$ 满足 $|\boldsymbol a+\boldsymbol b|=|\boldsymbol a-\boldsymbol b|$，则下面结论正确的是\mbox{（　　）}

（A）$\boldsymbol a\parallel\boldsymbol b$；（B）$\boldsymbol a\perp\boldsymbol b$；（C）$|\boldsymbol a|=|\boldsymbol b|$；（D）$\boldsymbol a+\boldsymbol b=\boldsymbol a-\boldsymbol b$。

\Needspace{4\baselineskip}
\textbf{解析}

\Needspace{5\baselineskip}
把已知等式两边平方：
\nopagebreak[4]
\[|\boldsymbol a+\boldsymbol b|^2=|\boldsymbol a-\boldsymbol b|^2.\]
\Needspace{5\baselineskip}
展开得
\nopagebreak[4]
\[\boldsymbol a^2+2\boldsymbol a\cdot\boldsymbol b+\boldsymbol b^2
=\boldsymbol a^2-2\boldsymbol a\cdot\boldsymbol b+\boldsymbol b^2,\]
\Needspace{5\baselineskip}
所以
\nopagebreak[4]
\[4\boldsymbol a\cdot\boldsymbol b=0,\]
\Needspace{5\baselineskip}
即
\nopagebreak[4]
\[\boldsymbol a\cdot\boldsymbol b=0.\]
因为 $\boldsymbol a$，$\boldsymbol b$ 都是非零向量，故 $\boldsymbol a\perp\boldsymbol b$。

选 B。

\gkanswer{答案：B。两边平方相减得 $4\boldsymbol a\cdot\boldsymbol b=0$，故 $\boldsymbol a\perp\boldsymbol b$。}

\end{answer}
\begin{answer}
\gksolhead{4}{}

\textbf{题目}：

已知命题 $p$：$\forall x_1,x_2\in\mathbf R$，$(f(x_2)-f(x_1))(x_2-x_1)\geqslant 0$，则 $\neg p$ 是\mbox{（　　）}

（A）$\exists x_1,x_2\in\mathbf R$，$(f(x_2)-f(x_1))(x_2-x_1)\leqslant 0$；（B）$\forall x_1,x_2\in\mathbf R$，$(f(x_2)-f(x_1))(x_2-x_1)\leqslant 0$；（C）$\exists x_1,x_2\in\mathbf R$，$(f(x_2)-f(x_1))(x_2-x_1)<0$；（D）$\forall x_1,x_2\in\mathbf R$，$(f(x_2)-f(x_1))(x_2-x_1)<0$。

\Needspace{4\baselineskip}
\textbf{解析}

\Needspace{5\baselineskip}
命题 $p$ 是全称命题，其否定要把“$\forall$”改为“$\exists$”，并把结论中的“$\geqslant$”改为“$<$”：
\nopagebreak[4]
\[\neg p:\ \exists x_1,x_2\in\mathbf R,\ (f(x_2)-f(x_1))(x_2-x_1)<0.\]

选 C。

\gkanswer{答案：C。全称命题的否定是存在命题，且不等号反向，得 $\exists x_1,x_2\in\mathbf R$，$(f(x_2)-f(x_1))(x_2-x_1)<0$。}

\end{answer}
\begin{answer}
\gksolhead{5}{}

\textbf{题目}：

一排 9 个座位坐了 3 个三口之家，若每家人坐在一起，则不同的坐法种数为\mbox{（　　）}

（A）$3\times3!$；（B）$3\times(3!)^3$；（C）$(3!)^4$；（D）$9!$。

\Needspace{4\baselineskip}
\textbf{解析}

\Needspace{5\baselineskip}
把每个三口之家看成一个整体，三个整体在 9 个座位上占三段连续座位，整体之间的排列有
\nopagebreak[4]
\[3!\]
种。

\Needspace{5\baselineskip}
每个家庭内部的 3 个人在自家这一段座位上有 $3!$ 种坐法，三个家庭共有
\nopagebreak[4]
\[(3!)^3\]
种内部排列。

\Needspace{5\baselineskip}
因此总坐法数为
\nopagebreak[4]
\[3!\cdot(3!)^3=(3!)^4=6^4=1296.\]

选 C。

\gkanswer{答案：C。三个家庭整体排列有 $3!$ 种，每家内部排列各有 $3!$ 种，共 $(3!)^4=1296$ 种。}

\end{answer}
\begin{answer}
\gksolhead{6}{}

\textbf{题目}：

在等差数列 $\{a_n\}$ 中，已知 $a_4+a_8=16$，则该数列前 11 项和 $S_{11}=$\mbox{（　　）}

（A）$58$；（B）$88$；（C）$143$；（D）$176$。

\Needspace{4\baselineskip}
\textbf{解析}

\Needspace{5\baselineskip}
等差数列中下标之和相等则对应项之和相等：
\nopagebreak[4]
\[a_1+a_{11}=a_4+a_8=16.\]
\Needspace{5\baselineskip}
前 11 项和为
\nopagebreak[4]
\[S_{11}=\frac{11(a_1+a_{11})}{2}=\frac{11\times16}{2}=88.\]

选 B。

\gkanswer{答案：B。$a_1+a_{11}=a_4+a_8=16$，故 $S_{11}=\dfrac{11(a_1+a_{11})}{2}=88$。}

\end{answer}
\begin{answer}
\gksolhead{7}{}

\textbf{题目}：

已知 $\sin\alpha-\cos\alpha=\sqrt2$，$\alpha\in(0,\pi)$，则 $\tan\alpha=$\mbox{（　　）}

（A）$-1$；（B）$-\dfrac{\sqrt2}{2}$；（C）$\dfrac{\sqrt2}{2}$；（D）$1$。

\Needspace{4\baselineskip}
\textbf{解析}

\Needspace{5\baselineskip}
将 $\sin\alpha-\cos\alpha=\sqrt2$ 两边平方：
\nopagebreak[4]
\[\sin^2\alpha-2\sin\alpha\cos\alpha+\cos^2\alpha=2.\]
\Needspace{5\baselineskip}
由 $\sin^2\alpha+\cos^2\alpha=1$ 得
\nopagebreak[4]
\[1-\sin2\alpha=2,\]
\Needspace{5\baselineskip}
即
\nopagebreak[4]
\[\sin2\alpha=-1.\]
\Needspace{5\baselineskip}
因为 $\alpha\in(0,\pi)$，所以 $2\alpha\in(0,2\pi)$。在此区间内正弦值为 $-1$ 的角只有
\nopagebreak[4]
\[2\alpha=\frac{3\pi}{2},\]
\Needspace{5\baselineskip}
故
\nopagebreak[4]
\[\alpha=\frac{3\pi}{4},\qquad \tan\alpha=-1.\]

选 A。

\gkanswer{答案：A。平方得 $\sin2\alpha=-1$，结合 $\alpha\in(0,\pi)$ 得 $\alpha=\dfrac{3\pi}{4}$，故 $\tan\alpha=-1$。}

\end{answer}
\begin{answer}
\gksolhead{8}{}

\textbf{题目}：

设变量 $x$，$y$ 满足 $\begin{cases}x-y\leqslant 10,\\0\leqslant x+y\leqslant 20,\\0\leqslant y\leqslant 15,\end{cases}$ 则 $2x+3y$ 的最大值为\mbox{（　　）}

（A）$20$；（B）$35$；（C）$45$；（D）$55$。

\Needspace{4\baselineskip}
\textbf{解析}

\Needspace{5\baselineskip}
先画可行域。边界为
\nopagebreak[4]
\[y=x-10,\qquad y=20-x,\qquad y=0,\qquad y=15.\]
\Needspace{5\baselineskip}
由 $y\geqslant x-10$、$y\leqslant20-x$、$0\leqslant y\leqslant15$ 得可行域是五边形，其顶点为
\nopagebreak[4]
\[(0,0),\ (10,0),\ (15,5),\ (5,15),\ (0,15).\]
\Needspace{5\baselineskip}
分别计算 $z=2x+3y$：
\nopagebreak[4]
\[\begin{aligned}
(0,0)&:\ 0,\\
(10,0)&:\ 20,\\
(15,5)&:\ 30+15=45,\\
(5,15)&:\ 10+45=55,\\
(0,15)&:\ 45.
\end{aligned}\]
最大值为 $55$。

选 D。

\gkanswer{答案：D。可行域顶点处目标值依次为 $0,20,45,55,45$，最大值为 $55$（在 $(5,15)$ 取得）。}

\end{answer}
\begin{answer}
\gksolhead{9}{}

\textbf{题目}：

执行如图所示的程序框图，则输出的 $S$ 值是\mbox{（　　）}
\nopagebreak[4]
\begin{center}% ============================================================
% f2012_liaoning_science_q09_flowchart.tex
% 2012 辽宁卷（理科）第 9 题：程序框图
% 来源 PDF：普通高考/2012/2012辽宁理.pdf
% 原页：第 1 页（A3），印面 751，第 9 题题干下方
% 初始 S=4、i=1；判断 i<9，是则 S=2/(2-S)、i=i+1 后循环，
% 否则输出 S 并结束。
% ============================================================
\begin{tikzpicture}[>=Stealth, line width=0.7pt, font=\small,
  every node/.style={inner sep=2pt},
  term/.style={draw, rounded corners=3pt, minimum width=2.0cm,
    minimum height=0.58cm, align=center},
  proc/.style={draw, minimum width=2.6cm, minimum height=0.58cm, align=center},
  io/.style={draw, trapezium, trapezium left angle=70, trapezium right angle=110,
    minimum width=2.4cm, minimum height=0.58cm, align=center},
  dec/.style={draw, diamond, aspect=2, minimum width=2.4cm,
    minimum height=0.9cm, align=center}]
  \node[term] (start) at (0,0) {开始};
  \node[proc] (initS) at (0,-0.92) {$S=4$};
  \node[proc] (initI) at (0,-1.84) {$i=1$};
  \node[dec] (test) at (0,-3.02) {$i<9$};
  \node[proc] (bodyS) at (0,-4.46) {$S=\dfrac{2}{2-S}$};
  \node[proc] (bodyI) at (0,-5.52) {$i=i+1$};
  \node[io] (out) at (3.35,-3.02) {输出 $S$};
  \node[term] (end) at (3.35,-4.22) {结束};

  \draw[->] (start) -- (initS);
  \draw[->] (initS) -- (initI);
  \draw[->] (initI) -- (test);
  \draw[->] (test.south) -- node[right] {是} (bodyS.north);
  \draw[->] (bodyS) -- (bodyI);
  % 循环体回到判断框左侧
  \draw[->] (bodyI.west) -- ++(-2.45,0) |- (test.west);
  % “否”分支输出
  \draw[->] (test.east) -- node[above] {否} (out.west);
  \draw[->] (out) -- (end);
\end{tikzpicture}
\end{center}
\nopagebreak[4]
（A）$-1$；（B）$\dfrac32$；（C）$\dfrac23$；（D）$4$。

\Needspace{4\baselineskip}
\textbf{解析}

按框图逐步计算。初始 $S=4$，$i=1$。

\Needspace{5\baselineskip}
当 $i<9$ 时执行
\nopagebreak[4]
\[S\leftarrow\frac{2}{2-S},\qquad i\leftarrow i+1.\]
\Needspace{5\baselineskip}
依次得
\nopagebreak[4]
\[\begin{aligned}
i=1&:\ S=\frac{2}{2-4}=-1,\\
i=2&:\ S=\frac{2}{2-(-1)}=\frac23,\\
i=3&:\ S=\frac{2}{2-\frac23}=\frac32,\\
i=4&:\ S=\frac{2}{2-\frac32}=4.
\end{aligned}\]
\Needspace{5\baselineskip}
此后 $S$ 以 $4,-1,\dfrac23,\dfrac32$ 为周期循环。执行到 $i=8$ 时又回到 $S=4$，随后 $i=9$，判断 $i<9$ 不成立，输出
\nopagebreak[4]
\[S=4.\]

选 D。

\gkanswer{答案：D。$S$ 以 $4,-\!1,\dfrac23,\dfrac32$ 为周期循环，$i=1,2,\ldots,8$ 共执行 8 次，$i=9$ 时输出 $S=4$。}

\end{answer}
\begin{answer}
\gksolhead{10}{}

\textbf{题目}：

在长为 $12\ \mathrm{cm}$ 的线段 $AB$ 上任取一点 $C$。现作一矩形，邻边长分别等于线段 $AC$，$CB$ 的长，则该矩形面积小于 $32\ \mathrm{cm}^2$ 的概率为\mbox{（　　）}

（A）$\dfrac16$；（B）$\dfrac13$；（C）$\dfrac23$；（D）$\dfrac45$。

\Needspace{4\baselineskip}
\textbf{解析}

\Needspace{5\baselineskip}
设 $AC=x$，则 $CB=12-x$，其中 $0<x<12$。矩形面积为
\nopagebreak[4]
\[S=x(12-x).\]
\Needspace{5\baselineskip}
条件 $S<32$ 即
\nopagebreak[4]
\[x(12-x)<32,\]
\Needspace{5\baselineskip}
整理得
\nopagebreak[4]
\[x^2-12x+32>0,\]
\Needspace{5\baselineskip}
即
\nopagebreak[4]
\[(x-4)(x-8)>0.\]
\Needspace{5\baselineskip}
所以
\nopagebreak[4]
\[0<x<4\quad\text{\text{或}}\quad 8<x<12.\]
\Needspace{5\baselineskip}
有利长度为
\nopagebreak[4]
\[4+(12-8)=8,\]
\Needspace{5\baselineskip}
故所求概率为
\nopagebreak[4]
\[P=\frac8{12}=\frac23.\]

选 C。

\gkanswer{答案：C。设 $AC=x$，则 $x(12-x)<32$，即 $x<4$ 或 $x>8$，有利长度 $8$，概率 $\dfrac8{12}=\dfrac23$。}

\end{answer}
\begin{answer}
\gksolhead{11}{}

\textbf{题目}：

设函数 $f(x)(x\in\mathbf R)$ 满足 $f(-x)=f(x)$，$f(x)=f(2-x)$，且当 $x\in[0,1]$ 时，$f(x)=x^3$。又函数 $g(x)=|x\cos(\pi x)|$，则函数 $h(x)=g(x)-f(x)$ 在 $\left[-\dfrac12,\dfrac32\right]$ 上的零点个数为\mbox{（　　）}

（A）$5$；（B）$6$；（C）$7$；（D）$8$。

\Needspace{4\baselineskip}
\textbf{解析}

先由对称性与周期性写出 $f(x)$ 在 $\left[-\dfrac12,\dfrac32\right]$ 上的表达式。

由 $f(-x)=f(x)$，$f$ 为偶函数；由 $f(x)=f(2-x)$，$f$ 关于 $x=1$ 对称，且周期为 $2$。

\Needspace{5\baselineskip}
当 $0\leqslant x\leqslant1$ 时，$f(x)=x^3$。因此
\nopagebreak[4]
\[f(x)=\begin{cases}
-x^3,&-\dfrac12\leqslant x<0,\\[2mm]
x^3,&0\leqslant x\leqslant1,\\[2mm]
(2-x)^3,&1<x\leqslant\dfrac32.
\end{cases}\]

零点方程为 $g(x)=f(x)$。

\Needspace{5\baselineskip}
\textbf{（i）$-\dfrac12\leqslant x\leqslant0$。} 此时 $x<0$，$\cos\pi x\geqslant0$，故
\nopagebreak[4]
\[g(x)=-x\cos\pi x,\qquad f(x)=-x^3.\]
\Needspace{5\baselineskip}
当 $x<0$ 时约去 $-x$，得
\nopagebreak[4]
\[\cos\pi x=x^2.\]
\Needspace{5\baselineskip}
令 $\varphi(x)=\cos\pi x-x^2$。则
\nopagebreak[4]
\[\varphi'(x)=-\pi\sin\pi x-2x>0\qquad\left(-\frac12<x<0\right),\]
\Needspace{5\baselineskip}
因为 $\sin\pi x<0$，$-2x>0$。又
\nopagebreak[4]
\[\varphi\left(-\frac12\right)=-\frac14<0,\qquad \varphi(0)=1>0,\]
故在 $\left(-\dfrac12,0\right)$ 内恰有一个零点。加上 $x=0$ 本身也是零点，本段共 \textbf{2 个}。

\Needspace{5\baselineskip}
\textbf{（ii）$0<x\leqslant\dfrac12$。} 此时 $\cos\pi x\geqslant0$，$g(x)=x\cos\pi x$，方程化为
\nopagebreak[4]
\[\cos\pi x=x^2.\]
\Needspace{5\baselineskip}
令 $\varphi(x)=\cos\pi x-x^2$，则
\nopagebreak[4]
\[\varphi'(x)=-\pi\sin\pi x-2x<0,\]
\Needspace{5\baselineskip}
且
\nopagebreak[4]
\[\varphi(0)=1>0,\qquad \varphi\left(\frac12\right)=-\frac14<0,\]
故本段恰有 \textbf{1 个}零点。

\Needspace{5\baselineskip}
\textbf{（iii）$\dfrac12<x\leqslant1$。} 此时 $\cos\pi x\leqslant0$，$g(x)=-x\cos\pi x$，方程化为
\nopagebreak[4]
\[-\cos\pi x=x^2.\]
\Needspace{5\baselineskip}
令
\nopagebreak[4]
\[\psi(x)=-\cos\pi x-x^2.\]
\Needspace{5\baselineskip}
则
\nopagebreak[4]
\[\psi''(x)=\pi^2\cos\pi x-2<0\qquad\left(\frac12<x<1\right),\]
\Needspace{5\baselineskip}
所以 $\psi$ 严格凹。又
\nopagebreak[4]
\[\psi\left(\frac12\right)=-\frac14<0,\qquad
\psi\left(\frac34\right)=\frac{\sqrt2}{2}-\frac9{16}>0,\qquad
\psi(1)=0.\]
严格凹函数至多有两个零点；结合上述符号，本段在 $\left(\dfrac12,1\right)$ 内有一个零点，再加上端点 $x=1$，共 \textbf{2 个}。

\Needspace{5\baselineskip}
\textbf{（iv）$1<x\leqslant\dfrac32$。} 此时 $\cos\pi x\leqslant0$，$g(x)=-x\cos\pi x$，方程为
\nopagebreak[4]
\[-x\cos\pi x=(2-x)^3.\]
\Needspace{5\baselineskip}
令
\nopagebreak[4]
\[\psi(x)=-x\cos\pi x-(2-x)^3.\]
\Needspace{5\baselineskip}
在 $\left(1,\dfrac32\right)$ 上，
\nopagebreak[4]
\[\psi''(x)=2\pi\sin\pi x+\pi^2x\cos\pi x-6(2-x)<0,\]
\Needspace{5\baselineskip}
故 $\psi$ 严格凹。又
\nopagebreak[4]
\[\psi(1)=0,\qquad \psi'(1)=4>0,\qquad
\psi\left(\frac32\right)=-\frac18<0,\]
所以 $\psi$ 在 $x=1$ 右侧先增后减，且恰有一个零点落在 $\left(1,\dfrac32\right)$ 内，本段新增 \textbf{1 个}。

\Needspace{5\baselineskip}
合计
\nopagebreak[4]
\[2+1+2+1=6.\]

选 B。

\gkanswer{答案：B。分四段讨论 $|\cos\pi x|=x^2$ 型方程，零点依次落在 $\left[-\dfrac12,0\right]$ 两个、$\left(0,\dfrac12\right)$ 一个、$\left(\dfrac12,1\right]$ 两个、$\left(1,\dfrac32\right]$ 一个，共 $6$ 个。}

\end{answer}
\begin{answer}
\gksolhead{12}{}

\textbf{题目}：

若 $x\in[0,+\infty)$，则下列不等式恒成立的是\mbox{（　　）}

（A）$\mathrm e^x\leqslant 1+x+x^2$；（B）$\dfrac1{\sqrt{1+x}}\leqslant 1-\dfrac12x+\dfrac14x^2$；（C）$\cos x\geqslant 1-\dfrac12x^2$；（D）$\ln(1+x)\geqslant x-\dfrac18x^2$。

\Needspace{4\baselineskip}
\textbf{解析}

逐项判断。

（A）取 $x=2$，则 $\mathrm e^2\approx7.389>1+2+4=7$，不恒成立。

\Needspace{5\baselineskip}
（B）取 $x=\dfrac12$，左边
\nopagebreak[4]
\[\frac1{\sqrt{1+\frac12}}=\sqrt{\frac23}\approx0.8165,\]
\Needspace{5\baselineskip}
右边
\nopagebreak[4]
\[1-\frac14+\frac1{16}=\frac{13}{16}=0.8125,\]
左边大于右边，不恒成立。

\Needspace{5\baselineskip}
（C）令
\nopagebreak[4]
\[\varphi(x)=\cos x-1+\frac12x^2.\]
\Needspace{5\baselineskip}
则
\nopagebreak[4]
\[\varphi'(x)=-\sin x+x=x-\sin x\geqslant0\qquad(x\geqslant0),\]
\Needspace{5\baselineskip}
因为 $x\geqslant\sin x$。又 $\varphi(0)=0$，所以 $\varphi(x)\geqslant0$，即
\nopagebreak[4]
\[\cos x\geqslant1-\frac12x^2\]
恒成立。

（D）取 $x=1$，左边 $\ln2\approx0.693$，右边 $1-\dfrac18=0.875$，不恒成立。

选 C。

\gkanswer{答案：C。令 $\varphi(x)=\cos x-1+\dfrac12x^2$，$\varphi'(x)=x-\sin x\geqslant0$，$\varphi(0)=0$，故（C）恒成立；其余三式分别取 $x=2,1,\dfrac12$ 可否定。}

\end{answer}
\begin{answer}
\gksolhead{13}{}

\textbf{题目}：

一个几何体的三视图如图所示，则该几何体的表面积为\gkblank{}。
\nopagebreak[4]
\begin{center}% ============================================================
% f2012_liaoning_science_q13_three_views.tex
% 2012 辽宁卷（理科）第 13 题：三视图
% 来源 PDF：普通高考/2012/2012辽宁理.pdf
% 原页：第 1 页（A3），印面 751，第 13 题题干下方
% 主视图 4×1，内部距左端 1、3 处有虚线；
% 侧视图 3×1，内部距左端 0.5、2.5 处有虚线；
% 俯视图 4×3 矩形，内含直径 2 的圆。
% ============================================================
\begin{tikzpicture}[>=Stealth, line width=0.7pt, font=\small,
  x=1.5cm, y=1.5cm,
  dim/.style={line width=0.4pt},
  every node/.style={inner sep=1.5pt}]
  % ---------- 主视图 ----------
  \draw (0,4) rectangle (4,5);
  \draw[dashed] (1,4) -- (1,5);
  \draw[dashed] (3,4) -- (3,5);
  % 主视图尺寸
  \draw[dim,<->] (0,3.55) -- (1,3.55) node[midway,below,font=\tiny,fill=white,inner sep=1pt] {$1$};
  \draw[dim,<->] (1,3.55) -- (3,3.55) node[midway,below,font=\tiny,fill=white,inner sep=1pt] {$2$};
  \draw[dim,<->] (3,3.55) -- (4,3.55) node[midway,below,font=\tiny,fill=white,inner sep=1pt] {$1$};

  % ---------- 侧视图 ----------
  \draw (5.3,4) rectangle (8.3,5);
  \draw[dashed] (5.8,4) -- (5.8,5);
  \draw[dashed] (7.8,4) -- (7.8,5);
  \draw[dim,<->] (5.3,3.55) -- (5.8,3.55) node[midway,below,font=\tiny,fill=white,inner sep=1pt] {$0.5$};
  \draw[dim,<->] (5.8,3.55) -- (7.8,3.55) node[midway,below,font=\tiny,fill=white,inner sep=1pt] {$2$};
  \draw[dim,<->] (7.8,3.55) -- (8.3,3.55) node[midway,below,font=\tiny,fill=white,inner sep=1pt] {$0.5$};
  % 侧视图高度
  \draw[dim,<->] (8.55,4) -- (8.55,5) node[midway,right,font=\tiny,fill=white,inner sep=1pt] {$1$};

  % ---------- 俯视图 ----------
  \draw (0,0) rectangle (4,3);
  \draw (2,1.5) circle (1);
\end{tikzpicture}
\end{center}
\nopagebreak[4]
\Needspace{4\baselineskip}
\textbf{解析}

由三视图还原几何体：

\begin{itemize}
\item 主视图外轮廓为 $4\times1$ 矩形，内部在距左端 $1$、$3$ 处有两条虚线；
\item 侧视图外轮廓为 $3\times1$ 矩形，内部在距左端 $0.5$、$2.5$ 处有两条虚线；
\item 俯视图为 $4\times3$ 矩形，内含一个直径为 $2$ 的圆。
\end{itemize}

因此几何体是一个 $4\times3\times1$ 的长方体，沿竖直方向挖去一个直径为 $2$ 的圆柱形通孔。

\Needspace{5\baselineskip}
长方体外表面积为
\nopagebreak[4]
\[S_0=2(4\times3+4\times1+3\times1)=2(12+4+3)=38.\]
\Needspace{5\baselineskip}
挖孔后，上、下底面各少一个半径为 $1$ 的圆，面积共减少
\nopagebreak[4]
\[2\pi\cdot1^2=2\pi;\]
\Needspace{5\baselineskip}
孔壁增加侧面积
\nopagebreak[4]
\[2\pi\cdot1\cdot1=2\pi.\]
\Needspace{5\baselineskip}
两者相抵，故总表面积为
\nopagebreak[4]
\[S=S_0-2\pi+2\pi=38.\]

填 $38$。

\gkanswer{答案：$38$。几何体是 $4\times3\times1$ 长方体挖去一个直径 $2$ 的圆柱形通孔，外表面积 $2(12+4+3)=38$，孔壁面积与上下两个圆孔面积相抵。}

\end{answer}
\begin{answer}
\gksolhead{14}{}

\textbf{题目}：

已知等比数列 $\{a_n\}$ 为递增数列，且 $a_5^2=a_{10}$，$2(a_n+a_{n+2})=5a_{n+1}$，则数列 $\{a_n\}$ 的通项公式 $a_n=$\gkblank{}。

\Needspace{4\baselineskip}
\textbf{解析}

\Needspace{5\baselineskip}
设公比为 $q$。由
\nopagebreak[4]
\[2(a_n+a_{n+2})=5a_{n+1}\]
\Needspace{5\baselineskip}
两边除以 $a_n$（$a_n>0$），得
\nopagebreak[4]
\[2(1+q^2)=5q,\]
\Needspace{5\baselineskip}
即
\nopagebreak[4]
\[2q^2-5q+2=0.\]
\Needspace{5\baselineskip}
解得
\nopagebreak[4]
\[q=2\quad\text{\text{或}}\quad q=\frac12.\]
\Needspace{5\baselineskip}
因为是递增数列且各项为正，故 $q>1$，所以
\nopagebreak[4]
\[q=2.\]

\Needspace{5\baselineskip}
再由
\nopagebreak[4]
\[a_5^2=a_{10}\]
\Needspace{5\baselineskip}
得
\nopagebreak[4]
\[(a_1q^4)^2=a_1q^9.\]
\Needspace{5\baselineskip}
因为 $a_1>0$，约去 $a_1q^8$，得
\nopagebreak[4]
\[a_1=q=2.\]
\Needspace{5\baselineskip}
所以
\nopagebreak[4]
\[a_n=2\cdot2^{n-1}=2^n.\]

填 $2^n$。

\gkanswer{答案：$2^n$。由 $2(1+q^2)=5q$ 得 $q=2$ 或 $\dfrac12$，递增取 $q=2$；再由 $a_5^2=a_{10}$ 得 $a_1=2$。}

\end{answer}
\begin{answer}
\gksolhead{15}{}

\textbf{题目}：

已知 $P$，$Q$ 为抛物线 $x^2=2y$ 上两点，点 $P$，$Q$ 的横坐标分别为 $4$，$-2$，过 $P$，$Q$ 分别作抛物线的切线，两切线交于 $A$，则点 $A$ 的纵坐标为\gkblank{}。

\Needspace{4\baselineskip}
\textbf{解析}

\Needspace{5\baselineskip}
抛物线为
\nopagebreak[4]
\[y=\frac12x^2,\qquad y'=x.\]

\Needspace{5\baselineskip}
点 $P$ 的横坐标为 $4$，故 $P(4,8)$，切线斜率为 $4$，切线方程为
\nopagebreak[4]
\[y-8=4(x-4),\]
\Needspace{5\baselineskip}
即
\nopagebreak[4]
\[y=4x-8.\]

\Needspace{5\baselineskip}
点 $Q$ 的横坐标为 $-2$，故 $Q(-2,2)$，切线斜率为 $-2$，切线方程为
\nopagebreak[4]
\[y-2=-2(x+2),\]
\Needspace{5\baselineskip}
即
\nopagebreak[4]
\[y=-2x-2.\]

\Needspace{5\baselineskip}
联立两切线：
\nopagebreak[4]
\[4x-8=-2x-2,\]
\Needspace{5\baselineskip}
得
\nopagebreak[4]
\[6x=6,\qquad x=1.\]
\Needspace{5\baselineskip}
代入得
\nopagebreak[4]
\[y=4-8=-4.\]
所以 $A$ 的纵坐标为 $-4$。

填 $-4$。

\gkanswer{答案：$-4$。两切线为 $y=4x-8$ 与 $y=-2x-2$，交点 $A(1,-4)$。}

\end{answer}
\begin{answer}
\gksolhead{16}{}

\textbf{题目}：

已知正三棱锥 $P-ABC$，点 $P$，$A$，$B$，$C$ 都在半径为 $\sqrt3$ 的球面上。若 $PA$，$PB$，$PC$ 两两互相垂直，则球心到截面 $ABC$ 的距离为\gkblank{}。

\Needspace{4\baselineskip}
\textbf{解析}

因为 $PA$，$PB$，$PC$ 两两垂直，以 $P$ 为原点，$PA$，$PB$，$PC$ 所在直线为三条坐标轴建立空间直角坐标系。

\Needspace{5\baselineskip}
设
\nopagebreak[4]
\[PA=PB=PC=l.\]
\Needspace{5\baselineskip}
则
\nopagebreak[4]
\[A(l,0,0),\qquad B(0,l,0),\qquad C(0,0,l).\]
\Needspace{5\baselineskip}
四面体 $PABC$ 的外接球球心为
\nopagebreak[4]
\[O\left(\frac l2,\frac l2,\frac l2\right),\]
\Needspace{5\baselineskip}
半径为
\nopagebreak[4]
\[R=\sqrt{3\left(\frac l2\right)^2}=\frac{\sqrt3}{2}l.\]
\Needspace{5\baselineskip}
由题意 $R=\sqrt3$，故
\nopagebreak[4]
\[\frac{\sqrt3}{2}l=\sqrt3,\]
\Needspace{5\baselineskip}
得
\nopagebreak[4]
\[l=2.\]

\Needspace{5\baselineskip}
截面 $ABC$ 的方程为
\nopagebreak[4]
\[x+y+z=2.\]
\Needspace{5\baselineskip}
球心 $O(1,1,1)$ 到该平面的距离为
\nopagebreak[4]
\[d=\frac{|1+1+1-2|}{\sqrt{1^2+1^2+1^2}}
=\frac1{\sqrt3}=\frac{\sqrt3}{3}.\]

填 $\dfrac{\sqrt3}{3}$。

\gkanswer{答案：$\dfrac{\sqrt3}{3}$。设 $PA=PB=PC=l$，外接球半径 $\dfrac{\sqrt3}{2}l=\sqrt3$ 得 $l=2$；球心到平面 $x+y+z=2$ 的距离为 $\dfrac{|3-2|}{\sqrt3}=\dfrac{\sqrt3}{3}$。}

\end{answer}
\begin{answer}
\gksolhead{17}{}

\textbf{题目}：

在 $\triangle ABC$ 中，角 $A$，$B$，$C$ 的对边分别为 $a$，$b$，$c$。角 $A$，$B$，$C$ 成等差数列。

（1） 求 $\cos B$ 的值；

（2） 若边 $a$，$b$，$c$ 成等比数列，求 $\sin A\sin C$ 的值。

\Needspace{4\baselineskip}
\textbf{解析}

\gksubpart{（1）}

\Needspace{5\baselineskip}
因为 $A$，$B$，$C$ 成等差数列，所以
\nopagebreak[4]
\[2B=A+C.\]
\Needspace{5\baselineskip}
又
\nopagebreak[4]
\[A+B+C=\pi,\]
\Needspace{5\baselineskip}
故
\nopagebreak[4]
\[3B=\pi,\]
\Needspace{5\baselineskip}
即
\nopagebreak[4]
\[B=\frac\pi3.\]
\Needspace{5\baselineskip}
所以
\nopagebreak[4]
\[\cos B=\cos\frac\pi3=\frac12.\]

\gksubpart{（2）}

\Needspace{5\baselineskip}
因为 $a$，$b$，$c$ 成等比数列，所以
\nopagebreak[4]
\[b^2=ac.\]
\Needspace{5\baselineskip}
由正弦定理
\nopagebreak[4]
\[a=2R\sin A,\qquad b=2R\sin B,\qquad c=2R\sin C,\]
\Needspace{5\baselineskip}
代入得
\nopagebreak[4]
\[(2R\sin B)^2=(2R\sin A)(2R\sin C),\]
\Needspace{5\baselineskip}
即
\nopagebreak[4]
\[\sin^2B=\sin A\sin C.\]
\Needspace{5\baselineskip}
由（1）知 $B=\dfrac\pi3$，所以
\nopagebreak[4]
\[\sin A\sin C=\sin^2\frac\pi3=\left(\frac{\sqrt3}{2}\right)^2=\frac34.\]

\gkanswer{答案（1）：$\dfrac12$。$2B=A+C$ 且 $A+B+C=\pi$，故 $B=\dfrac\pi3$，$\cos B=\dfrac12$。}

\gkanswer{答案（2）：$\dfrac34$。由 $b^2=ac$ 及正弦定理得 $\sin^2B=\sin A\sin C$，故 $\sin A\sin C=\dfrac34$。}

\end{answer}
\begin{answer}
\gksolhead{18}{}

\textbf{题目}：

如图，直三棱柱 $ABC-A'B'C'$，$\angle BAC=90^\circ$，$AB=AC=\lambda AA'$，点 $M$，$N$ 分别为 $A'B$ 和 $B'C'$ 的中点。
\nopagebreak[4]
\begin{center}% ============================================================
% f2012_liaoning_science_q18_prism.tex
% 2012 辽宁卷（理科）第 18 题：直三棱柱与二面角
% 来源 PDF：普通高考/2012/2012辽宁理.pdf
% 原页：第 1 页（A3），印面 751，第 18 题题干右侧
% 直三棱柱 ABC-A'B'C'，M 为 A'B 中点，N 为 B'C' 中点。
% 虚线表示被遮挡的棱 AB、AC、AA'。
% ============================================================
\begin{tikzpicture}[>=Stealth, line width=0.7pt, font=\small,
  every node/.style={inner sep=1.5pt}]
  \coordinate (A)  at (1.00,1.60);
  \coordinate (B)  at (0.20,0.20);
  \coordinate (C)  at (3.20,0.80);
  \coordinate (Ap) at (1.00,3.40);
  \coordinate (Bp) at (0.20,2.00);
  \coordinate (Cp) at (3.20,2.60);
  \coordinate (M)  at (0.60,1.80);
  \coordinate (N)  at (1.70,2.30);

  % 被遮挡的棱
  \draw[dashed] (A) -- (B);
  \draw[dashed] (A) -- (C);
  \draw[dashed] (A) -- (Ap);

  % 可见棱
  \draw (B) -- (C);
  \draw (B) -- (Bp);
  \draw (C) -- (Cp);
  \draw (Ap) -- (Bp);
  \draw (Bp) -- (Cp);
  \draw (Cp) -- (Ap);
  \draw (Ap) -- (B);

  % 中点连线与二面角相关线段
  \draw (M) -- (N);
  \draw (Ap) -- (N);
  \draw (C) -- (M);
  \draw (C) -- (N);

  % 点与标签
  \fill (A) circle (1.1pt) node[below right] {$A$};
  \fill (B) circle (1.1pt) node[below] {$B$};
  \fill (C) circle (1.1pt) node[right] {$C$};
  \fill (Ap) circle (1.1pt) node[above] {$A'$};
  \fill (Bp) circle (1.1pt) node[left] {$B'$};
  \fill (Cp) circle (1.1pt) node[right] {$C'$};
  \fill (M) circle (1.0pt) node[below left] {$M$};
  \fill (N) circle (1.0pt) node[above right] {$N$};
\end{tikzpicture}
\end{center}
\nopagebreak[4]
（1） 证明：$MN\parallel$ 平面 $A'ACC'$；

（2） 若二面角 $A'-MN-C$ 为直二面角，求 $\lambda$ 的值。

\Needspace{4\baselineskip}
\textbf{解析}

\Needspace{5\baselineskip}
设 $AA'=1$，则 $AB=AC=\lambda$。以 $A$ 为原点，$AB$，$AC$，$AA'$ 所在直线分别为 $x$，$y$，$z$ 轴建立空间直角坐标系，则
\nopagebreak[4]
\[\begin{aligned}
&A(0,0,0),\quad B(\lambda,0,0),\quad C(0,\lambda,0),\\
&A'(0,0,1),\quad B'(\lambda,0,1),\quad C'(0,\lambda,1).
\end{aligned}\]

\gksubpart{（1）}

\Needspace{5\baselineskip}
$M$ 为 $A'B$ 的中点，$N$ 为 $B'C'$ 的中点，故
\nopagebreak[4]
\[M\left(\frac\lambda2,0,\frac12\right),\qquad
N\left(\frac\lambda2,\frac\lambda2,1\right).\]
\Needspace{5\baselineskip}
于是
\nopagebreak[4]
\[\overrightarrow{MN}=\left(0,\frac\lambda2,\frac12\right)
=\frac12(0,\lambda,1).\]
\Needspace{5\baselineskip}
而
\nopagebreak[4]
\[\overrightarrow{AC'}=(0,\lambda,1).\]
\Needspace{5\baselineskip}
所以
\nopagebreak[4]
\[\overrightarrow{MN}=\frac12\overrightarrow{AC'},\]
即 $MN\parallel AC'$。

\Needspace{5\baselineskip}
因为 $AC'\subset$ 平面 $A'ACC'$，且 $M\notin$ 平面 $A'ACC'$，所以
\nopagebreak[4]
\[MN\parallel\text{\text{平面} }A'ACC'.\]

\gksubpart{（2）}

\Needspace{5\baselineskip}
平面 $A'MN$ 内有两个向量
\nopagebreak[4]
\[\overrightarrow{A'M}=\left(\frac\lambda2,0,-\frac12\right),\qquad
\overrightarrow{A'N}=\left(\frac\lambda2,\frac\lambda2,0\right).\]
\Needspace{5\baselineskip}
取法向量
\nopagebreak[4]
\[\boldsymbol n_1=\overrightarrow{A'M}\times\overrightarrow{A'N}
=\left(\frac\lambda4,-\frac\lambda4,\frac{\lambda^2}{4}\right)
\sim(1,-1,\lambda).\]

\Needspace{5\baselineskip}
平面 $CMN$ 内有两个向量
\nopagebreak[4]
\[\overrightarrow{CM}=\left(\frac\lambda2,-\lambda,\frac12\right),\qquad
\overrightarrow{CN}=\left(\frac\lambda2,-\frac\lambda2,1\right).\]
\Needspace{5\baselineskip}
取法向量
\nopagebreak[4]
\[\boldsymbol n_2=\overrightarrow{CM}\times\overrightarrow{CN}
=\left(-\frac{3\lambda}{4},-\frac{\lambda}{4},\frac{\lambda^2}{4}\right)
\sim(-3,-1,\lambda).\]

\Needspace{5\baselineskip}
二面角 $A'-MN-C$ 为直二面角，等价于两平面的法向量垂直：
\nopagebreak[4]
\[\boldsymbol n_1\cdot\boldsymbol n_2=0.\]
\Needspace{5\baselineskip}
即
\nopagebreak[4]
\[1\cdot(-3)+(-1)\cdot(-1)+\lambda\cdot\lambda=0,\]
\Needspace{5\baselineskip}
化简得
\nopagebreak[4]
\[\lambda^2-2=0.\]
\Needspace{5\baselineskip}
因为 $\lambda>0$，所以
\nopagebreak[4]
\[\lambda=\sqrt2.\]

\gkanswer{答案（1）：见解析版}

\gkanswer{答案（2）：$\sqrt2$。取 $AA'=1$，两半平面的法向量分别为 $(1,-1,\lambda)$、$(-3,-1,\lambda)$，由法向量垂直得 $\lambda^2=2$。}

\end{answer}
\begin{answer}
\gksolhead{19}{}

\textbf{题目}：

电视传媒公司为了解某地区电视观众对某类体育节目的收视情况，随机抽取了 100 名观众进行调查。下面是根据调查结果绘制的观众日均收看该体育节目时间的频率分布直方图：
\nopagebreak[4]
\begin{center}% ============================================================
% f2012_liaoning_science_q19_histogram.tex
% 2012 辽宁卷（理科）第 19 题：频率分布直方图
% 来源 PDF：普通高考/2012/2012辽宁理.pdf
% 原页：第 2 页（A3），印面 752，第 19 题题干下方
% 各组频率/组距依次为 0.010、0.018、0.020、0.025、0.022、0.005，
% 组距为 10 分钟，共 6 组。
% ============================================================
\begin{tikzpicture}[>=Stealth, line width=0.7pt, font=\small,
  x=1.2cm, y=1cm,
  every node/.style={inner sep=1.5pt}]

  % 坐标轴
  \draw[->] (-0.25,0) -- (7.75,0) node[right] {分钟};
  \draw[->] (0,-0.25) -- (0,3.75) node[above] {$\dfrac{\text{频率}}{\text{组距}}$};

  % 各矩形
  \draw (0,0) rectangle (1.2,1.30);   % 0.010
  \draw (1.2,0) rectangle (2.4,2.34); % 0.018
  \draw (2.4,0) rectangle (3.6,2.60); % 0.020
  \draw (3.6,0) rectangle (4.8,3.25); % 0.025
  \draw (4.8,0) rectangle (6.0,2.86); % 0.022
  \draw (6.0,0) rectangle (7.2,0.65); % 0.005

  % 高度虚线
  \draw[dashed] (0,1.30) -- (1.2,1.30);
  \draw[dashed] (0,2.34) -- (2.4,2.34);
  \draw[dashed] (0,2.60) -- (3.6,2.60);
  \draw[dashed] (0,3.25) -- (4.8,3.25);
  \draw[dashed] (0,2.86) -- (6.0,2.86);
  \draw[dashed] (0,0.65) -- (7.2,0.65);

  % 纵轴刻度值
  \node[left,font=\scriptsize] at (0,0.65) {$0.005$};
  \node[left,font=\scriptsize] at (0,1.30) {$0.010$};
  \node[left,font=\scriptsize] at (0,2.34) {$0.018$};
  \node[left,font=\scriptsize] at (0,2.60) {$0.020$};
  \node[left,font=\scriptsize] at (0,2.86) {$0.022$};
  \node[left,font=\scriptsize] at (0,3.25) {$0.025$};

  % 横轴刻度值
  \foreach \i/\t in {0/0,1/10,2/20,3/30,4/40,5/50,6/60}
    \node[below,font=\scriptsize] at (\i*1.2,0) {$\t$};
\end{tikzpicture}
\end{center}
\nopagebreak[4]
将日均收看该体育节目时间不低于 40 分钟的观众称为“体育迷”。

（1） 根据已知条件完成下面的 $2\times2$ 列联表，并据此资料你是否认为“体育迷”与性别有关？

\begin{center}
\begin{tabular}{llll}
\toprule
 & 非体育迷 & 体育迷 & 合计 \\
\midrule
男 &  &  &  \\
女 &  & $10$ & $55$ \\
合计 &  &  &  \\
\bottomrule
\end{tabular}
\end{center}

（2） 将上述调查所得到的频率视为概率。现在从该地区大量电视观众中，采用随机抽样方法每次抽取一名观众，抽取 3 次，记被抽取的 3 名观众中的“体育迷”人数为 $X$。若每次抽取的结果是相互独立的，求 $X$ 的分布列、期望 $E(X)$ 和方差 $D(X)$。

\Needspace{10\baselineskip}
附：$\chi^2=\dfrac{n(n_{11}n_{22}-n_{12}n_{21})^2}{n_{1+}n_{2+}n_{+1}n_{+2}}$，

\begin{center}
\begin{tabular}{lll}
\toprule
$P(\chi^2\geqslant k)$ & $0.05$ & $0.01$ \\
\midrule
$k$ & $3.841$ & $6.635$ \\
\bottomrule
\end{tabular}
\end{center}

\Needspace{4\baselineskip}
\textbf{解析}

\gksubpart{（1）}

\Needspace{5\baselineskip}
由频率分布直方图，各组频率为“组距 $10$ 乘以 $\dfrac{\text{\text{频率}}}{\text{\text{组距}}}$”，依次为
\nopagebreak[4]
\[0.10,\ 0.18,\ 0.20,\ 0.25,\ 0.22,\ 0.05.\]
\Needspace{5\baselineskip}
样本容量为 $100$，故各组人数依次为
\nopagebreak[4]
\[10,\ 18,\ 20,\ 25,\ 22,\ 5.\]
\Needspace{5\baselineskip}
日均收看时间不低于 $40$ 分钟的人数为
\nopagebreak[4]
\[22+5=27,\]
即“体育迷”共 $27$ 人。

\Needspace{5\baselineskip}
女性共 $55$ 人，其中“体育迷”$10$ 人，故女性非体育迷
\nopagebreak[4]
\[55-10=45\text{ \text{人}}.\]
\Needspace{5\baselineskip}
男性“体育迷”为
\nopagebreak[4]
\[27-10=17\text{ \text{人}},\]
\Needspace{5\baselineskip}
男性共
\nopagebreak[4]
\[100-55=45\text{ \text{人}},\]
\Needspace{5\baselineskip}
男性非体育迷为
\nopagebreak[4]
\[45-17=28\text{ \text{人}}.\]
列联表为

\begin{center}
\begin{tabular}{llll}
\toprule
 & 非体育迷 & 体育迷 & 合计 \\
\midrule
男 & $28$ & $17$ & $45$ \\
女 & $45$ & $10$ & $55$ \\
合计 & $73$ & $27$ & $100$ \\
\bottomrule
\end{tabular}
\end{center}

\Needspace{5\baselineskip}
计算
\nopagebreak[4]
\[\begin{aligned}
\chi^2
&=\frac{n(n_{11}n_{22}-n_{12}n_{21})^2}
{n_{1+}n_{2+}n_{+1}n_{+2}}\\
&=\frac{100(28\times10-17\times45)^2}
{45\times55\times73\times27}\\
&=\frac{100\times(-485)^2}{4878225}\\
&=\frac{940900}{195129}\approx4.822.
\end{aligned}\]
\Needspace{5\baselineskip}
因为
\nopagebreak[4]
\[4.822>3.841,\]
所以有 $95\%$ 的把握认为“体育迷”与性别有关。

\gksubpart{（2）}

\Needspace{5\baselineskip}
把频率视为概率，任取一名观众是“体育迷”的概率为
\nopagebreak[4]
\[p=\frac{27}{100}.\]
\Needspace{5\baselineskip}
抽取 3 次，每次相互独立，故
\nopagebreak[4]
\[X\sim B\left(3,\frac{27}{100}\right).\]
\Needspace{5\baselineskip}
分布列为
\nopagebreak[4]
\[\begin{aligned}
P(X=0)&=\mathrm C_3^0\left(\frac{27}{100}\right)^0\left(\frac{73}{100}\right)^3
=\frac{389017}{1000000},\\
P(X=1)&=\mathrm C_3^1\left(\frac{27}{100}\right)\left(\frac{73}{100}\right)^2
=\frac{431649}{1000000},\\
P(X=2)&=\mathrm C_3^2\left(\frac{27}{100}\right)^2\left(\frac{73}{100}\right)
=\frac{159651}{1000000},\\
P(X=3)&=\mathrm C_3^3\left(\frac{27}{100}\right)^3
=\frac{19683}{1000000}.
\end{aligned}\]
\Needspace{5\baselineskip}
所以
\nopagebreak[4]
\[E(X)=np=3\times\frac{27}{100}=\frac{81}{100},\]
\nopagebreak[4]
\[D(X)=np(1-p)=3\times\frac{27}{100}\times\frac{73}{100}
=\frac{5913}{10000}.\]

\gkanswer{答案（1）：列联表为男 $28,17,45$，女 $45,10,55$，合计 $73,27,100$；$\chi^2=\dfrac{940900}{195129}\approx4.822>3.841$，有 $95\%$ 的把握认为“体育迷”与性别有关。}

\gkanswer{答案（2）：$X\sim B\left(3,\dfrac{27}{100}\right)$，分布列为 $P(X=0)=\dfrac{389017}{1000000}$，$P(X=1)=\dfrac{431649}{1000000}$，$P(X=2)=\dfrac{159651}{1000000}$，$P(X=3)=\dfrac{19683}{1000000}$；$E(X)=\dfrac{81}{100}$，$D(X)=\dfrac{5913}{10000}$。}

\end{answer}
\begin{answer}
\gksolhead{20}{}

\textbf{题目}：

如图，椭圆 $C_0:\dfrac{x^2}{a^2}+\dfrac{y^2}{b^2}=1$（$a>b>0$，$a$，$b$ 为常数），动圆 $C_1:x^2+y^2=t_1^2$，$b<t_1<a$。点 $A_1$，$A_2$ 分别为 $C_0$ 的左、右顶点，$C_1$ 与 $C_0$ 相交于 $A$，$B$，$C$，$D$ 四点。
\nopagebreak[4]
\begin{center}% ============================================================
% f2012_liaoning_science_q20_ellipse.tex
% 2012 辽宁卷（理科）第 20 题：椭圆与动圆
% 来源 PDF：普通高考/2012/2012辽宁理.pdf
% 原页：第 2 页（A3），印面 752，第 20 题题干下方
% 椭圆 C0: x^2/a^2+y^2/b^2=1，动圆 C1: x^2+y^2=t^2。
% 图中取 a=3、b=2、t=2.3 作示意，交点记为 A、B、C、D。
% ============================================================
\begin{tikzpicture}[>=Stealth, line width=0.7pt, font=\small,
  x=1.1cm, y=1.1cm,
  every node/.style={inner sep=1.5pt}]

  % 坐标轴
  \draw[->] (-3.8,0) -- (3.9,0) node[right] {$x$};
  \draw[->] (0,-2.7) -- (0,2.8) node[above] {$y$};
  \node[below left] at (0,0) {$O$};

  % 椭圆与动圆
  \draw (0,0) ellipse (3 and 2);
  \draw (0,0) circle (2.3);

  % 四个交点（a=3,b=2,t=2.3 时 x≈1.524,y≈1.722）
  \coordinate (A) at (-1.524,1.722);
  \coordinate (B) at (-1.524,-1.722);
  \coordinate (C) at (1.524,-1.722);
  \coordinate (D) at (1.524,1.722);

  \fill (A) circle (1.1pt) node[above left] {$A$};
  \fill (B) circle (1.1pt) node[below left] {$B$};
  \fill (C) circle (1.1pt) node[below right] {$C$};
  \fill (D) circle (1.1pt) node[above right] {$D$};

  % 左右顶点
  \fill (-3,0) circle (1.1pt) node[below left] {$A_1$};
  \fill (3,0) circle (1.1pt) node[below right] {$A_2$};
\end{tikzpicture}
\end{center}
\nopagebreak[4]
（1） 求直线 $AA_1$ 与直线 $A_2B$ 交点 $M$ 的轨迹方程；

（2） 设动圆 $C_2:x^2+y^2=t_2^2$ 与 $C_0$ 相交于 $A'$，$B'$，$C'$，$D'$ 四点，其中 $b<t_2<a$，$t_1\neq t_2$。若矩形 $ABCD$ 与矩形 $A'B'C'D'$ 的面积相等，证明：$t_1^2+t_2^2$ 为定值。

\Needspace{4\baselineskip}
\textbf{解析}

\gksubpart{（1）}

\Needspace{5\baselineskip}
设
\nopagebreak[4]
\[A(-x_1,y_1),\qquad B(-x_1,-y_1),\qquad C(x_1,-y_1),\qquad D(x_1,y_1),\]
\Needspace{5\baselineskip}
其中 $x_1>0$，$y_1>0$。则 $A$，$B$，$C$，$D$ 在圆 $C_1$ 上，故
\nopagebreak[4]
\[x_1^2+y_1^2=t_1^2.\]
\Needspace{5\baselineskip}
又它们在椭圆 $C_0$ 上，故
\nopagebreak[4]
\[\frac{x_1^2}{a^2}+\frac{y_1^2}{b^2}=1.\]

\Needspace{5\baselineskip}
椭圆左、右顶点为
\nopagebreak[4]
\[A_1(-a,0),\qquad A_2(a,0).\]

\Needspace{5\baselineskip}
直线 $AA_1$ 的斜率为
\nopagebreak[4]
\[k_1=\frac{y_1}{a-x_1},\]
\Needspace{5\baselineskip}
方程为
\nopagebreak[4]
\[y=\frac{y_1}{a-x_1}(x+a).\]
\Needspace{5\baselineskip}
直线 $A_2B$ 的斜率为
\nopagebreak[4]
\[k_2=\frac{-y_1}{-x_1-a}=\frac{y_1}{a+x_1},\]
\Needspace{5\baselineskip}
方程为
\nopagebreak[4]
\[y=\frac{y_1}{a+x_1}(x-a).\]

\Needspace{5\baselineskip}
设交点 $M(x,y)$，联立得
\nopagebreak[4]
\[\frac{x+a}{a-x_1}=\frac{x-a}{a+x_1}.\]
\Needspace{5\baselineskip}
去分母：
\nopagebreak[4]
\[(x+a)(a+x_1)=(x-a)(a-x_1).\]
\Needspace{5\baselineskip}
展开：
\nopagebreak[4]
\[ax+xx_1+a^2+ax_1=ax-xx_1-a^2+ax_1.\]
\Needspace{5\baselineskip}
消去 $ax$、$ax_1$，得
\nopagebreak[4]
\[2xx_1=-2a^2,\]
\Needspace{5\baselineskip}
即
\nopagebreak[4]
\[x=-\frac{a^2}{x_1}.\]
\Needspace{5\baselineskip}
代回直线 $AA_1$：
\nopagebreak[4]
\[\begin{aligned}
y&=\frac{y_1}{a-x_1}\left(-\frac{a^2}{x_1}+a\right)\\
&=\frac{y_1}{a-x_1}\cdot\frac{a(x_1-a)}{x_1}\\
&=-\frac{ay_1}{x_1}.
\end{aligned}\]
\Needspace{5\baselineskip}
于是
\nopagebreak[4]
\[x_1=-\frac{a^2}{x},\qquad y_1=\frac{ay}{x}.\]
\Needspace{5\baselineskip}
代入椭圆方程：
\nopagebreak[4]
\[\frac{1}{a^2}\cdot\frac{a^4}{x^2}
+\frac{1}{b^2}\cdot\frac{a^2y^2}{x^2}=1.\]
\Needspace{5\baselineskip}
即
\nopagebreak[4]
\[\frac{a^2}{x^2}+\frac{a^2y^2}{b^2x^2}=1.\]
\Needspace{5\baselineskip}
两边乘以 $\dfrac{x^2}{a^2}$，得
\nopagebreak[4]
\[\frac{x^2}{a^2}-\frac{y^2}{b^2}=1.\]
\Needspace{5\baselineskip}
由 $x=-\dfrac{a^2}{x_1}$ 且 $0<x_1<a$，知 $x<-a$，所以轨迹为双曲线
\nopagebreak[4]
\[\frac{x^2}{a^2}-\frac{y^2}{b^2}=1\]
的左支。

\gksubpart{（2）}

\Needspace{5\baselineskip}
对动圆半径 $t$（$b<t<a$），设交点仍为
\nopagebreak[4]
\[A(-x,y),\quad B(-x,-y),\quad C(x,-y),\quad D(x,y),\]
\Needspace{5\baselineskip}
其中 $x>0$，$y>0$。则
\nopagebreak[4]
\[x^2+y^2=t^2,\qquad \frac{x^2}{a^2}+\frac{y^2}{b^2}=1.\]
\Needspace{5\baselineskip}
由第二式与 $y^2=t^2-x^2$ 得
\nopagebreak[4]
\[\frac{x^2}{a^2}+\frac{t^2-x^2}{b^2}=1.\]
\Needspace{5\baselineskip}
整理：
\nopagebreak[4]
\[x^2\left(\frac1{a^2}-\frac1{b^2}\right)=1-\frac{t^2}{b^2},\]
\Needspace{5\baselineskip}
即
\nopagebreak[4]
\[x^2=\frac{a^2(t^2-b^2)}{a^2-b^2}.\]
\Needspace{5\baselineskip}
同理
\nopagebreak[4]
\[y^2=\frac{b^2(a^2-t^2)}{a^2-b^2}.\]
\Needspace{5\baselineskip}
矩形 $ABCD$ 的面积为
\nopagebreak[4]
\[S=2x\cdot2y=4xy,\]
\Needspace{5\baselineskip}
故
\nopagebreak[4]
\[S^2=16x^2y^2
=\frac{16a^2b^2(t^2-b^2)(a^2-t^2)}{(a^2-b^2)^2}.\]

\Needspace{5\baselineskip}
设两矩形面积相等。因为 $t_1\neq t_2$，且面积均为正，故
\nopagebreak[4]
\[(t_1^2-b^2)(a^2-t_1^2)
=(t_2^2-b^2)(a^2-t_2^2).\]
\Needspace{5\baselineskip}
令 $u_1=t_1^2$，$u_2=t_2^2$，展开：
\nopagebreak[4]
\[(a^2+b^2)u_1-u_1^2-a^2b^2
=(a^2+b^2)u_2-u_2^2-a^2b^2.\]
\Needspace{5\baselineskip}
移项并因式分解：
\nopagebreak[4]
\[(u_1-u_2)\bigl[(a^2+b^2)-(u_1+u_2)\bigr]=0.\]
\Needspace{5\baselineskip}
因为 $t_1\neq t_2$，故 $u_1\neq u_2$，于是
\nopagebreak[4]
\[u_1+u_2=a^2+b^2,\]
\Needspace{5\baselineskip}
即
\nopagebreak[4]
\[t_1^2+t_2^2=a^2+b^2.\]
所以 $t_1^2+t_2^2$ 为定值 $a^2+b^2$。

\gkanswer{答案（1）：$\dfrac{x^2}{a^2}-\dfrac{y^2}{b^2}=1$（$x<-a$）。设 $A(-x_1,y_1)$，$B(-x_1,-y_1)$，联立两直线得 $M\left(-\dfrac{a^2}{x_1},-\dfrac{ay_1}{x_1}\right)$，消去 $x_1,y_1$ 即得。}

\gkanswer{答案（2）：见解析版}

\end{answer}
\begin{answer}
\gksolhead{21}{}

\textbf{题目}：

设 $f(x)=\ln(x+1)+\sqrt{x+1}+ax+b$（$a$，$b\in\mathbf R$，$a$，$b$ 为常数），曲线 $y=f(x)$ 与直线 $y=\dfrac32x$ 在 $(0,0)$ 点相切。

（1） 求 $a$，$b$ 的值；

（2） 证明：当 $0<x<2$ 时，$f(x)<\dfrac{9x}{x+6}$。

\Needspace{4\baselineskip}
\textbf{解析}

\gksubpart{（1）}

\Needspace{5\baselineskip}
曲线 $y=f(x)$ 与直线 $y=\dfrac32x$ 在 $(0,0)$ 点相切，包含两个条件：
\nopagebreak[4]
\[f(0)=0,\qquad f'(0)=\frac32.\]

\Needspace{5\baselineskip}
先由 $f(0)=0$：
\nopagebreak[4]
\[f(0)=\ln1+\sqrt1+b=1+b=0,\]
\Needspace{5\baselineskip}
得
\nopagebreak[4]
\[b=-1.\]

\Needspace{5\baselineskip}
再求导：
\nopagebreak[4]
\[f'(x)=\frac1{x+1}+\frac1{2\sqrt{x+1}}+a.\]
\Needspace{5\baselineskip}
于是
\nopagebreak[4]
\[f'(0)=1+\frac12+a=\frac32,\]
\Needspace{5\baselineskip}
得
\nopagebreak[4]
\[a=0.\]

\Needspace{5\baselineskip}
所以
\nopagebreak[4]
\[a=0,\qquad b=-1.\]

\gksubpart{（2）}

\Needspace{5\baselineskip}
此时
\nopagebreak[4]
\[f(x)=\ln(x+1)+\sqrt{x+1}-1.\]
\Needspace{5\baselineskip}
令
\nopagebreak[4]
\[t=\sqrt{x+1}.\]
\Needspace{5\baselineskip}
当 $0<x<2$ 时，
\nopagebreak[4]
\[1<t<\sqrt3<2,\]
\Needspace{5\baselineskip}
且
\nopagebreak[4]
\[x=t^2-1.\]
\Needspace{5\baselineskip}
于是
\nopagebreak[4]
\[\begin{aligned}
f(x)&=\ln t^2+t-1\\
&=2\ln t+t-1.
\end{aligned}\]
\Needspace{5\baselineskip}
而
\nopagebreak[4]
\[\frac{9x}{x+6}
=\frac{9(t^2-1)}{t^2+5}
=\frac{9(t-1)(t+1)}{t^2+5}.\]

\Needspace{5\baselineskip}
由熟知不等式
\nopagebreak[4]
\[\ln t\leqslant t-1\qquad(t>0),\]
\Needspace{5\baselineskip}
得
\nopagebreak[4]
\[2\ln t+t-1\leqslant2(t-1)+t-1=3(t-1).\]
\Needspace{5\baselineskip}
因为 $t>1$，所以 $t-1>0$。又 $1<t<2$，有
\nopagebreak[4]
\[t^2-3t+2=(t-1)(t-2)<0,\]
\Needspace{5\baselineskip}
即
\nopagebreak[4]
\[t^2+5<3t+3=3(t+1).\]
\Needspace{5\baselineskip}
两边同乘 $\dfrac{t-1}{t^2+5}>0$，得
\nopagebreak[4]
\[3(t-1)<\frac{9(t-1)(t+1)}{t^2+5}.\]
\Needspace{5\baselineskip}
因此
\nopagebreak[4]
\[f(x)=2\ln t+t-1\leqslant3(t-1)
<\frac{9(t-1)(t+1)}{t^2+5}
=\frac{9x}{x+6}.\]
\Needspace{5\baselineskip}
故当 $0<x<2$ 时，
\nopagebreak[4]
\[f(x)<\frac{9x}{x+6}.\]

\gkanswer{答案（1）：$a=0$，$b=-1$。由 $f(0)=0$ 得 $b=-1$，由 $f'(0)=\dfrac32$ 得 $1+\dfrac12+a=\dfrac32$，故 $a=0$。}

\gkanswer{答案（2）：见解析版}

\end{answer}
\begin{answer}
\gksolhead{22}{}

\textbf{题目}：

【A】如图，$\odot O$ 和 $\odot O'$ 相交于 $A$，$B$ 两点，过 $A$ 作两圆的切线分别交两圆于 $C$，$D$ 两点，连接 $DB$ 并延长交 $\odot O$ 于点 $E$。证明：
\nopagebreak[4]
\begin{center}% ============================================================
% f2012_liaoning_science_q22_circles.tex
% 2012 辽宁卷（理科）第 22 题：两圆相交与切线
% 来源 PDF：普通高考/2012/2012辽宁理.pdf
% 原页：第 2 页（A3），印面 752，第 22 题题干下方
% 圆 O 与圆 O' 交于 A、B；过 A 的两条切线分别交两圆于 C、D；
% DB 延长交圆 O 于 E。
% 取 O(-1.2,0)、O'(1.3,0)，A(0,1.5)、B(0,-1.5)，
% 两圆半径分别为 sqrt(3.69)、sqrt(3.94)。
% ============================================================
\begin{tikzpicture}[>=Stealth, line width=0.7pt, font=\small,
  x=1.1cm, y=1.1cm,
  every node/.style={inner sep=1.5pt}]
  \coordinate (O)  at (-1.2,0);
  \coordinate (Op) at (1.3,0);
  \coordinate (A)  at (0,1.5);
  \coordinate (B)  at (0,-1.5);
  \coordinate (C)  at (-2.856,-0.975);
  \coordinate (D)  at (3.049,-0.939);
  \coordinate (E)  at (-1.788,-1.829);

  % 两圆
  \draw (O) circle ({sqrt(3.69)});
  \draw (Op) circle ({sqrt(3.94)});

  % 切线、弦与割线
  \draw (A) -- (C);
  \draw (A) -- (D);
  \draw (A) -- (B);
  \draw (D) -- (B);
  \draw (B) -- (E);
  \draw (A) -- (E);
  \draw (B) -- (C);

  % 圆心
  \fill (O) circle (1.0pt) node[below right] {$O$};
  \fill (Op) circle (1.0pt) node[below left] {$O'$};

  % 点与标签
  \fill (A) circle (1.1pt) node[above] {$A$};
  \fill (B) circle (1.1pt) node[below] {$B$};
  \fill (C) circle (1.1pt) node[below left] {$C$};
  \fill (D) circle (1.1pt) node[below right] {$D$};
  \fill (E) circle (1.1pt) node[left] {$E$};
\end{tikzpicture}
\end{center}
\nopagebreak[4]
（1） $AC\cdot BD=AD\cdot AB$；

（2） $AC=AE$。

\Needspace{4\baselineskip}
\textbf{解析}

\gksubpart{（1）}

\Needspace{5\baselineskip}
因为 $AC$ 是 $\odot O'$ 在 $A$ 处的切线，$AB$ 是 $\odot O'$ 的弦，由弦切角定理，
\nopagebreak[4]
\[\angle CAB=\angle ADB.\]
\Needspace{5\baselineskip}
因为 $AD$ 是 $\odot O$ 在 $A$ 处的切线，$AB$ 是 $\odot O$ 的弦，同理
\nopagebreak[4]
\[\angle ACB=\angle DAB.\]
\Needspace{5\baselineskip}
于是
\nopagebreak[4]
\[\triangle ACB\sim\triangle DAB.\]
\Needspace{5\baselineskip}
对应边成比例：
\nopagebreak[4]
\[\frac{AC}{DA}=\frac{AB}{DB}.\]
\Needspace{5\baselineskip}
交叉相乘得
\nopagebreak[4]
\[AC\cdot BD=AD\cdot AB.\]

\gksubpart{（2）}

\Needspace{5\baselineskip}
在 $\triangle ABD$ 中，
\nopagebreak[4]
\[\angle ABD=\pi-\angle DAB-\angle ADB.\]
\Needspace{5\baselineskip}
由（1）知
\nopagebreak[4]
\[\angle DAB=\angle ACB,\qquad \angle ADB=\angle CAB.\]
\Needspace{5\baselineskip}
所以
\nopagebreak[4]
\[\begin{aligned}
\angle ABD
&=\pi-\angle ACB-\angle CAB\\
&=\angle ABC.
\end{aligned}\]
\Needspace{5\baselineskip}
因为 $E$，$B$，$D$ 共线，且 $B$ 在 $D$，$E$ 之间，所以 $BE$ 与 $BD$ 是反向延长线，故
\nopagebreak[4]
\[\angle ABE=\pi-\angle ABD=\pi-\angle ABC.\]

\Needspace{5\baselineskip}
四边形 $ACEB$ 内接于 $\odot O$，对角互补：
\nopagebreak[4]
\[\angle ACE+\angle ABE=\pi.\]
\Needspace{5\baselineskip}
于是
\nopagebreak[4]
\[\angle ACE=\pi-\angle ABE=\angle ABC.\]
\Needspace{5\baselineskip}
又在 $\odot O$ 中，同弧 $AC$ 所对的圆周角相等：
\nopagebreak[4]
\[\angle AEC=\angle ABC.\]
\Needspace{5\baselineskip}
因此
\nopagebreak[4]
\[\angle ACE=\angle AEC.\]
\Needspace{5\baselineskip}
由等角对等边，
\nopagebreak[4]
\[AC=AE.\]

\gkanswer{答案（1）：见解析版}

\gkanswer{答案（2）：见解析版}

\end{answer}
\begin{answer}
\gksolhead{23}{}

\textbf{题目}：

【B】在直角坐标系 $xOy$ 中，圆 $C_1:x^2+y^2=4$，圆 $C_2:(x-2)^2+y^2=4$。

（1） 在以 $O$ 为极点，$x$ 轴非负半轴为极轴的极坐标系中，分别写出圆 $C_1$，$C_2$ 的极坐标方程，并求出圆 $C_1$，$C_2$ 的交点坐标（用极坐标表示）；

（2） 求圆 $C_1$ 与 $C_2$ 的公共弦的参数方程。

\Needspace{4\baselineskip}
\textbf{解析}

\gksubpart{（1）}

\Needspace{5\baselineskip}
圆 $C_1:x^2+y^2=4$ 的极坐标方程为
\nopagebreak[4]
\[\rho^2=4,\]
\Needspace{5\baselineskip}
即
\nopagebreak[4]
\[\rho=2.\]

\Needspace{5\baselineskip}
圆 $C_2:(x-2)^2+y^2=4$ 展开为
\nopagebreak[4]
\[x^2+y^2-4x+4=4,\]
\Needspace{5\baselineskip}
即
\nopagebreak[4]
\[x^2+y^2=4x.\]
\Needspace{5\baselineskip}
代入 $x=\rho\cos\theta$，$x^2+y^2=\rho^2$，得
\nopagebreak[4]
\[\rho^2=4\rho\cos\theta.\]
\Needspace{5\baselineskip}
因为 $\rho=0$ 表示极点，已包含在方程中，故
\nopagebreak[4]
\[\rho=4\cos\theta.\]

\Needspace{5\baselineskip}
求交点，联立
\nopagebreak[4]
\[\rho=2,\qquad \rho=4\cos\theta.\]
\Needspace{5\baselineskip}
得
\nopagebreak[4]
\[2=4\cos\theta,\]
\Needspace{5\baselineskip}
即
\nopagebreak[4]
\[\cos\theta=\frac12.\]
\Needspace{5\baselineskip}
所以
\nopagebreak[4]
\[\theta=\pm\frac\pi3.\]
\Needspace{5\baselineskip}
交点极坐标为
\nopagebreak[4]
\[\left(2,\frac\pi3\right),\qquad \left(2,-\frac\pi3\right).\]

\gksubpart{（2）}

\Needspace{5\baselineskip}
两圆方程相减：
\nopagebreak[4]
\[[x^2+y^2]-[(x-2)^2+y^2]=4-4,\]
\Needspace{5\baselineskip}
即
\nopagebreak[4]
\[x^2-(x^2-4x+4)=0,\]
\Needspace{5\baselineskip}
得
\nopagebreak[4]
\[4x-4=0,\qquad x=1.\]
所以公共弦所在直线为 $x=1$。

\Needspace{5\baselineskip}
将 $x=1$ 代入 $x^2+y^2=4$：
\nopagebreak[4]
\[1+y^2=4,\qquad y^2=3,\]
\Needspace{5\baselineskip}
故
\nopagebreak[4]
\[-\sqrt3\leqslant y\leqslant\sqrt3.\]
\Needspace{5\baselineskip}
公共弦的参数方程为
\nopagebreak[4]
\[\begin{cases}
x=1,\\
y=t,
\end{cases}\qquad -\sqrt3\leqslant t\leqslant\sqrt3.\]

\gkanswer{答案（1）：$C_1:\rho=2$，$C_2:\rho=4\cos\theta$；交点极坐标为 $\left(2,\dfrac\pi3\right)$，$\left(2,-\dfrac\pi3\right)$。}

\gkanswer{答案（2）：$\begin{cases}x=1,\\y=t\end{cases}$（$-\sqrt3\leqslant t\leqslant\sqrt3$）。两圆方程相减得公共弦所在直线 $x=1$，再与 $x^2+y^2=4$ 联立。}

\end{answer}
\begin{answer}
\gksolhead{24}{}

\textbf{题目}：

【C】已知 $f(x)=|ax+1|$（$a\in\mathbf R$），不等式 $f(x)\leqslant 3$ 的解集为 $\{x\mid -2\leqslant x\leqslant 1\}$。

（1） 求 $a$ 的值；

（2） 若 $\left|f(x)-2f\left(\dfrac x2\right)\right|\leqslant k$ 恒成立，求 $k$ 的取值范围。

\Needspace{4\baselineskip}
\textbf{解析}

\gksubpart{（1）}

\Needspace{5\baselineskip}
不等式
\nopagebreak[4]
\[|ax+1|\leqslant3\]
\Needspace{5\baselineskip}
等价于
\nopagebreak[4]
\[-3\leqslant ax+1\leqslant3,\]
\Needspace{5\baselineskip}
即
\nopagebreak[4]
\[-4\leqslant ax\leqslant2.\]

\Needspace{5\baselineskip}
若 $a>0$，解集为
\nopagebreak[4]
\[\left[-\frac4a,\frac2a\right].\]
\Needspace{5\baselineskip}
与 $\{x\mid-2\leqslant x\leqslant1\}$ 比较，得
\nopagebreak[4]
\[-\frac4a=-2,\qquad \frac2a=1,\]
\Needspace{5\baselineskip}
两式都给出
\nopagebreak[4]
\[a=2.\]

\Needspace{5\baselineskip}
若 $a<0$，解集为
\nopagebreak[4]
\[\left[\frac2a,-\frac4a\right],\]
\Needspace{5\baselineskip}
需
\nopagebreak[4]
\[\frac2a=-2,\qquad -\frac4a=1,\]
第一式得 $a=-1$，第二式得 $a=-4$，矛盾。

若 $a=0$，不等式化为 $1\leqslant3$，解集为 $\mathbf R$，不符合。

\Needspace{5\baselineskip}
所以
\nopagebreak[4]
\[a=2.\]

\gksubpart{（2）}

\Needspace{5\baselineskip}
由（1），
\nopagebreak[4]
\[f(x)=|2x+1|.\]
\Needspace{5\baselineskip}
于是
\nopagebreak[4]
\[\begin{aligned}
\left|f(x)-2f\left(\frac x2\right)\right|
&=\left||2x+1|-2\left|\frac{2x}{2}+1\right|\right|\\
&=\left||2x+1|-2|x+1|\right|.
\end{aligned}\]
\Needspace{5\baselineskip}
令
\nopagebreak[4]
\[h(x)=|2x+1|-2|x+1|.\]
以 $x=-1$、$x=-\dfrac12$ 为分界点讨论。

\Needspace{5\baselineskip}
当 $x\leqslant-1$ 时，
\nopagebreak[4]
\[2x+1<0,\qquad x+1\leqslant0,\]
\Needspace{5\baselineskip}
故
\nopagebreak[4]
\[h(x)=-(2x+1)+2(x+1)=1.\]

\Needspace{5\baselineskip}
当 $-1\leqslant x\leqslant-\dfrac12$ 时，
\nopagebreak[4]
\[2x+1\leqslant0,\qquad x+1\geqslant0,\]
\Needspace{5\baselineskip}
故
\nopagebreak[4]
\[h(x)=-(2x+1)-2(x+1)=-4x-3.\]
\Needspace{5\baselineskip}
此时
\nopagebreak[4]
\[h(-1)=1,\qquad h\left(-\frac12\right)=-1,\]
\Needspace{5\baselineskip}
且 $h$ 单调递减，故
\nopagebreak[4]
\[-1\leqslant h(x)\leqslant1.\]

\Needspace{5\baselineskip}
当 $x\geqslant-\dfrac12$ 时，
\nopagebreak[4]
\[2x+1\geqslant0,\qquad x+1>0,\]
\Needspace{5\baselineskip}
故
\nopagebreak[4]
\[h(x)=(2x+1)-2(x+1)=-1.\]

\Needspace{5\baselineskip}
综上，
\nopagebreak[4]
\[-1\leqslant h(x)\leqslant1,\]
\Needspace{5\baselineskip}
即
\nopagebreak[4]
\[\left|f(x)-2f\left(\frac x2\right)\right|\leqslant1.\]
\Needspace{5\baselineskip}
要使不等式恒成立，必须且只需
\nopagebreak[4]
\[k\geqslant1.\]
\Needspace{5\baselineskip}
所以 $k$ 的取值范围是
\nopagebreak[4]
\[[1,+\infty).\]

\gkanswer{答案（1）：$a=2$。$|ax+1|\leqslant3$ 即 $-4\leqslant ax\leqslant2$，与 $-2\leqslant x\leqslant1$ 比较得 $a=2$。}

\gkanswer{答案（2）：$k\geqslant1$。令 $h(x)=|2x+1|-2|x+1|$，分段求得 $h(x)\in[-1,1]$，故 $k$ 的取值范围是 $[1,+\infty)$。}

\end{answer}